\documentclass[unnumsec,webpdf,contemporary,large]{oup-authoring-template}

\usepackage{microtype}
\usepackage{cleveref}

\crefname{figure}{Fig.}{Figs.}
\Crefname{figure}{Figure}{Figures}
\crefname{table}{Table}{Tables}
\Crefname{table}{Table}{Tables}
\crefname{section}{Section}{Sections}
\Crefname{section}{Section}{Sections}

\graphicspath{{figures/}}


\usepackage{tikz}
\usetikzlibrary{arrows.meta,positioning,backgrounds,fit,calc}

\usepackage{fontawesome5}
\hypersetup{colorlinks=true,urlcolor=blue,linkcolor=black,citecolor=black}

% Cross-reference the separately-compiled supplement (compile supplementary.tex
% first so supplementary.aux exists; labels are imported with a supp- prefix).
\usepackage{xr}
\externaldocument[supp-]{supplementary}

\begin{document}

\journaltitle{Bioinformatics}
\copyrightyear{2026}
\pubyear{2026}
\appnotes{Application Note}

\firstpage{1}

\title[RIsearch]{RIsearch and siOFF: An integrated, high-performance framework for RNA{--}RNA interaction and siRNA off-target assessment}

\author[1]{Stefano Roncelli \ORCID{0009-0001-4830-202X}}
\author[1]{Lorenzo Favaro \ORCID{0009-0002-8879-0555}}
\author[1]{Christian Anthon \ORCID{0000-0002-4336-1688}}
\author[1,$\ast$]{Jan Gorodkin \ORCID{0000-0001-5823-4000}}

\address[1]{\orgdiv{Center for non-coding RNA in Technology and Health, Section for Health Data Science and AI, Department of Public Health}, \orgname{University of Copenhagen}, \orgaddress{\street{Thorvaldsensvej 57}, \postcode{1871}, \state{Frederiksberg}, \country{Denmark}}}

\corresp[$\ast$]{Corresponding author. \href{mailto:gorodkin@rth.dk}{gorodkin@rth.dk}}

\received{Date}{0}{Year}
\revised{Date}{0}{Year}
\accepted{Date}{0}{Year}

\abstract{
\textbf{Motivation:} A substantial part of RNA function comes from its interactions with other RNAs.
Given the sheer scale of the search space (genome/transcriptome), there is high demand for fast RNA interaction search software that should also be user-friendly and easy to integrate with other tools.\\
\textbf{Results:} Here, we introduce new versions of RIsearch and its downstream siRNA off-target assessment pipeline.
Both tools now have substantially improved usability, interoperability, and installation procedures.
RIsearch~(v3) is a Rust reimplementation of the RIsearch2 search algorithm, verified against the reference; it reworks the original command-line interface, adds Python bindings, and improves performance.
siOFF modernises the previous siRNA off-target pipeline, through a redesigned CLI, improved distribution and faster execution.\\
\textbf{Availability and implementation:} RIsearch and siOFF are released under the Business Source License 1.1 at \href{https://github.com/saiden89/risearch}{github.com/saiden89/risearch} and \href{https://github.com/lorenzo-22/RIsearch_pipeline}{github.com/lorenzo-22/RIsearch\_pipeline}.
Both software packages are available on PyPI as \texttt{risearch} and \texttt{sioff}. The Rust binary is indexed in .\\
\textbf{Contact:} Jan Gorodkin, \href{mailto:gorodkin@rth.dk}{gorodkin@rth.dk}}

\keywords{RNA-RNA interactions, siRNA off-target assessment, Rust, Python}

\maketitle

\section{Introduction and motivation}

Intermolecular RNA interactions with both RNA and DNA underpin a wide range of regulatory mechanisms, and demand is increasing for predicting them at genome- or transcriptome-wide scale.
Whereas the RNA--RNA interactions continue with increasing attention, recent work has revealed enormous potential for RNA--DNA interactions~\citep{Lambolez2026RNA}.
When studying and using endogenous and synthetic small interfering RNAs (siRNAs), the lack of the off-target landscape can confound RNAi knockdown experiments and, for example, limit therapeutic application through unintended interactions with off-targets.

RIsearch2~\citep{alkan2017risearch2} is a tool designed for large-scale screens of near-complementary RNA--RNA interactions.
Because it ignores intramolecular interactions, it is particularly suitable for probing the off-target landscape of small RNAs, like siRNAs.
RIsearch combines a suffix-array index of the target with a seed-and-extend dynamic-programming step that uses a nearest-neighbour simplified energy model~\citep{wenzel2012risearch}.
The companion siRNA off-target pipeline turns these predictions into per-transcript binding probabilities and a transcriptome-wide off-targeting potential through a partition-function model that integrates transcript abundance and binding-site accessibility.

With the development of modern software ecosystems, new opportunities have emerged to improve how RIsearch2 and its pipeline were originally distributed and used.
Here, we introduce new versions of RIsearch2 and the siRNA pipeline, completely redesigned for improved usability, interoperability, installation, and development.
RIsearch~(v3) is a faster and memory-safe reimplementation of the same core search algorithm, while also providing a redesigned command-line interface, and PyO3 Python bindings.
The siRNA Off-Target pipeline, siOFF, is a rewrite of the previous siRNA off-target pipeline and can now use the RIsearch~(v3) bindings directly; it also serves as an example of using RIsearch embeddings and integrating them into a Python workflow.

\begin{figure*}[t]
\centering
\resizebox{\linewidth}{!}{% RIsearch + RIOT overview — TikZ source (Figure 1).
% A two-lane install -> execute -> output flow, blocks aligned on one row per lane.
%   Lane 1: RIsearch  (Crates.io / PyPI install, SA seed-and-extend, generic output).
%   Lane 2: RIOT      (PyPI install, RIsearch-backed pipeline, final off-target file).
% Required libraries (already loaded in main.tex for summary.tex):
%   \usetikzlibrary{arrows.meta,positioning,backgrounds,fit,calc}

% ── Palette (matches the package's two-tool colour identity) ─────────────────
\definecolor{CoreBlue}{RGB}{179,208,232}     % RIsearch fill
\definecolor{CoreBorder}{RGB}{68,114,196}    % RIsearch border
\definecolor{PipeGreen}{RGB}{198,224,180}    % RIOT fill
\definecolor{PipeBorder}{RGB}{83,153,51}     % RIOT border
\definecolor{CodeBg}{RGB}{246,246,248}       % code-snippet background
\definecolor{CodeBd}{RGB}{200,200,205}       % code-snippet border
\definecolor{ChanBg}{RGB}{255,255,255}       % install-channel chip
\definecolor{InkGray}{RGB}{70,70,70}
\definecolor{VrnaGold}{RGB}{229,178,72}      % ViennaRNA accent
\definecolor{NclsPlum}{RGB}{150,110,170}     % NCLS accent

\begin{tikzpicture}[
  font=\sffamily,
  >={Stealth[length=2.6mm]},
  laneband/.style={
    rounded corners=3pt, fill=#1, text=white, anchor=south west,
    minimum height=6mm, inner xsep=6pt, font=\sffamily\bfseries\small,
  },
  stage/.style={
    rounded corners=6pt, draw=#1, line width=0.9pt, fill=#1!18,
    align=center, inner sep=5pt, font=\sffamily\small\bfseries, text=InkGray,
  },
  info/.style={
    rounded corners=3pt, draw=#1!55, line width=0.5pt, fill=#1!8,
    align=left, inner sep=5pt, font=\sffamily\scriptsize, text=black,
  },
  code/.style={
    rounded corners=3pt, draw=CodeBd, line width=0.5pt, fill=CodeBg,
    align=left, inner sep=4pt, font=\ttfamily\scriptsize, text=black,
  },
  chip/.style={
    rounded corners=3pt, draw=#1, line width=0.7pt, fill=ChanBg,
    align=center, inner sep=3pt, font=\sffamily\scriptsize\bfseries, text=#1,
  },
  flow/.style={->, line width=1pt, draw=#1},
  caplbl/.style={font=\sffamily\scriptsize, text=InkGray, align=center},
]

% column centres (shared by both lanes so the two rows line up vertically)
\def\colA{0}
\def\colB{6.2}
\def\colC{12.4}

% =====================================================================
%  LANE 1 — RIsearch  (all blocks on one row, tops aligned at y = 0)
% =====================================================================
% ---- (1a) Install: two channels stacked vertically (no overlap) -----
\node[stage=CoreBorder, text width=3.4cm, anchor=north] (rs-install) at (\colA,0)
  {Install (Rust core)};
\node[chip=CoreBorder, anchor=north, text width=1.8cm]
  (rs-crates) at ([yshift=-4pt]rs-install.south) {Crates.io};
\node[code, anchor=north, text width=3.35cm]
  (rs-crates-cmd) at ([yshift=-2pt]rs-crates.south)
  {\$ cargo install risearch\\ \$ risearch --help};
\node[chip=CoreBorder, anchor=north, text width=1.8cm]
  (rs-pypi) at ([yshift=-5pt]rs-crates-cmd.south) {PyPI};
\node[code, anchor=north, text width=3.35cm]
  (rs-pypi-cmd) at ([yshift=-2pt]rs-pypi.south)
  {\$ pip install risearch\\ >>> import risearch};
\node[draw=CoreBorder!45, line width=0.6pt, rounded corners=6pt,
      fit=(rs-install)(rs-crates-cmd)(rs-pypi-cmd), inner sep=4pt] (rs-install-box) {};

% ---- (1b) Execute: conceptual description (no commands) -------------
\node[stage=CoreBorder, text width=4.55cm, anchor=north] (rs-run) at (\colB,0)
  {Index \& search};
\node[info=CoreBorder, anchor=north, text width=4.55cm]
  (rs-run-info) at ([yshift=-4pt]rs-run.south)
  {\textbullet~Suffix-array (SA) index of targets\\[2pt]
   \textbullet~Seed-and-extend DP extension\\[2pt]
   \textbullet~Simplified RIsearch energy model};
\node[draw=CoreBorder!45, line width=0.6pt, rounded corners=6pt,
      fit=(rs-run)(rs-run-info), inner sep=4pt] (rs-run-box) {};

% ---- (1c) Output: generic, illustrative (not the real schema) -------
\node[stage=CoreBorder, text width=5.0cm, anchor=north] (rs-out) at (\colC,0)
  {RIsearch output};
\node[code, anchor=north, text width=5.0cm]
  (rs-out-tbl) at ([yshift=-4pt]rs-out.south)
  {query   target   start   end   energy\\[1pt]
   q1      chrX      881    901   -28.4\\
   q1      chrY      233    253   -24.1\\[3pt]
   \# TSV output};
\node[draw=CoreBorder!45, line width=0.6pt, rounded corners=6pt,
      fit=(rs-out)(rs-out-tbl), inner sep=4pt] (rs-out-box) {};

% horizontal lane arrows (both pinned to the run-header centre line)
\draw[flow=CoreBorder] (rs-install-box.east|-rs-run) -- (rs-run-box.west|-rs-run);
\draw[flow=CoreBorder] (rs-run-box.east|-rs-run) -- (rs-out-box.west|-rs-run);

% lane-1 label tab
\node[laneband=CoreBorder]
  (rs-tab) at ([yshift=4pt]rs-install-box.north west)
  {RIsearch \textemdash\ RNA--RNA interaction search core};

% bounding helper for clearance below lane 1
\node[fit=(rs-install-box)(rs-run-box)(rs-out-box), inner sep=0] (lane1) {};

% =====================================================================
%  LANE 2 — RIOT  (all blocks on one row, tops aligned below lane 1)
% =====================================================================
\coordinate (l2) at ([yshift=-2.0cm]lane1.south -| rs-install);

% ---- (2a) Install ---------------------------------------------------
\node[stage=PipeBorder, text width=4.1cm, anchor=north] (ro-install) at (l2)
  {Install (RIOT pipeline)};
\node[chip=PipeBorder, anchor=north, text width=1.8cm]
  (ro-pypi) at ([yshift=-4pt]ro-install.south) {PyPI};
\node[code, anchor=north, text width=4.1cm]
  (ro-install-cmd) at ([yshift=-2pt]ro-pypi.south)
  {\$ pip install riot-rna\\[1pt]
   \# bundles: risearch wheel,\\
   \# ViennaRNA API, NCLS\\
   \# uv-locked, no binaries};
\node[draw=PipeBorder!45, line width=0.6pt, rounded corners=6pt,
      fit=(ro-install)(ro-install-cmd), inner sep=4pt] (ro-install-box) {};

% ---- (2b) Pipeline (built on RIsearch) ------------------------------
\node[stage=PipeBorder, text width=4.55cm, anchor=north]
  (ro-run) at ([xshift=\colB cm]ro-install.north) {Run off-target pipeline};
\node[code, anchor=north, text width=4.55cm]
  (ro-run-cmd) at ([yshift=-4pt]ro-run.south)
  {\$ riot off-targets \textbackslash\\
   \ \ \ -s sirna.fa --genome hg38.fa \textbackslash\\
   \ \ \ -t annot.gtf -a access/ \textbackslash\\
   \ \ \ --temperature 37 -o off.tsv\\[2pt]
   \# or: riot -c run.yaml};
\node[stage=CoreBorder, text width=4.4cm, anchor=north, font=\sffamily\scriptsize\bfseries]
  (st-search) at ([yshift=-4pt]ro-run-cmd.south) {1. RIsearch index + search (in-process)};
\node[stage=NclsPlum, text width=4.4cm, anchor=north, font=\sffamily\scriptsize\bfseries, fill=NclsPlum!18, draw=NclsPlum]
  (st-isect) at ([yshift=-3pt]st-search.south) {2. Intersect annotation \textbullet\ NCLS};
\node[stage=VrnaGold, text width=4.4cm, anchor=north, font=\sffamily\scriptsize\bfseries, fill=VrnaGold!22, draw=VrnaGold]
  (st-acc) at ([yshift=-3pt]st-isect.south) {3. Accessibility \textbullet\ ViennaRNA API};
\node[stage=PipeBorder, text width=4.4cm, anchor=north, font=\sffamily\scriptsize\bfseries]
  (st-prob) at ([yshift=-3pt]st-acc.south) {4. Boltzmann partition function $Z_s$};
\node[draw=PipeBorder!45, line width=0.6pt, rounded corners=6pt,
      fit=(ro-run)(ro-run-cmd)(st-search)(st-prob), inner sep=4pt] (ro-run-box) {};

\draw[flow=InkGray, line width=0.7pt] (st-search.south) -- (st-isect.north);
\draw[flow=InkGray, line width=0.7pt] (st-isect.south) -- (st-acc.north);
\draw[flow=InkGray, line width=0.7pt] (st-acc.south) -- (st-prob.north);

% ---- (2c) Final output file -----------------------------------------
\node[stage=PipeBorder, text width=5.0cm, anchor=north]
  (ro-out) at ([xshift=\colC cm]ro-install.north) {RIOT output};
\node[code, anchor=north, text width=5.0cm]
  (ro-out-tbl) at ([yshift=-4pt]ro-out.south)
  {sirna\_id   transcript\_id   P\_off\\[1pt]
   si\_001     ENST...0412      0.973\\
   si\_001     ENST...1097      0.118\\[3pt]
   \# TSV output};
\node[draw=PipeBorder!45, line width=0.6pt, rounded corners=6pt,
      fit=(ro-out)(ro-out-tbl), inner sep=4pt] (ro-out-box) {};

% horizontal lane arrows (both pinned to the run-header centre line)
\draw[flow=PipeBorder] (ro-install-box.east|-ro-run) -- (ro-run-box.west|-ro-run);
\draw[flow=PipeBorder] (ro-run-box.east|-ro-run) -- (ro-out-box.west|-ro-run);

% lane-2 label tab
\node[laneband=PipeBorder]
  (ro-tab) at ([yshift=4pt]ro-install-box.north west)
  {RIOT \textemdash\ RIsearch siRNA Off-Target pipeline};

% ---- "built on RIsearch" connector: straight segments, routed down the
%      clear corridor between the pipeline box and the off-target table so
%      it passes beside/under those boxes rather than through them.
\coordinate (cdrop) at ([yshift=-0.5cm]rs-out-box.south);            % just below Output
\coordinate (ccorr) at (9.2, 0 |- cdrop);                            % left into the corridor
\coordinate (cdown) at (9.2, 0 |- st-search.east);                   % straight down the corridor
\draw[flow=CoreBorder, line width=1pt, dashed]
  (rs-out-box.south) -- (cdrop) -- (ccorr) -- (cdown) -- (st-search.east);
\node[caplbl, fill=white, inner sep=1.5pt]
  at ([yshift=-0.3cm]$(cdrop)!0.5!(ccorr)$)
  {same engine, same process};

\end{tikzpicture}
}
\caption{
  \textbf{Installation and execution of RIsearch and the siOFF pipeline.}
  \textbf{(Top, blue)} \textbf{RIsearch} is a Rust reimplementation of the RIsearch2 search algorithm, installable via \texttt{cargo install risearch} (Crates.io) or as a Python wheel via \texttt{pip install risearch} (PyPI).
  The core builds a suffix-array index of target sequences and runs a seed-and-extend dynamic-programming search using the simplified RIsearch energy model, producing interaction hits as a TSV file.
  \textbf{(Bottom, green)} \textbf{siOFF} (siRNA Off-Target pipeline) is a timely Python~3 rewrite installable with a single \texttt{pip install sioff-rna}, which bundles the RIsearch dependencies, the ViennaRNA~2.7.2 Python API, and an in-memory nested containment list (NCLS) interval index with no external binaries required.
  siOFF orchestrates four stages: (\textit{i})~RIsearch index and search invoked natively via PyO3 bindings; (\textit{ii})~intersection of predicted binding sites with a transcriptome annotation using an NCLS index; (\textit{iii})~per-site RNA accessibility computed through the ViennaRNA Python API; and (\textit{iv})~a per-siRNA Boltzmann partition function yielding off-target probabilities, written to a TSV output file.
  The dashed arrow indicates that siOFF calls RIsearch within the same process via the Python bindings, with no subprocess or intermediate files.
}
\label{fig:summary}
\end{figure*}

A limitation of the original RIsearch2 package is how it is distributed and installed, as it is provided as C source code that each user must compile locally.
RIsearch~(v3) improves on these aspects while meeting three essential requirements at once:
\begin{enumerate}
\item retaining the genome-scale throughput that keeps whole-transcriptome screens practical;
\item installing from a standard package registry, with no manual build;
\item providing  a Python library, so that any analysis can  depend on it directly.
\end{enumerate}
C reaches that throughput but has no standard package registry of its own, so prebuilt distribution depends on external packaging, and it cannot guarantee memory safety at a library boundary.
Python distributes and embeds easily, but pure Python executes the seed-search and dynamic-programming loops in an interpreter, which is far slower than compiled native code, and so cannot reach that throughput.
Rust, however, meets all three: it compiles to fast native code, installs from standard package registries, and can exposes the algorithm through a memory-safe Python library.
These capabilities are why we rebuilt RIsearch in Rust.

\section{RIsearch reimplementation}\label{sec:core}\label{subsec:pybindings}
RIsearch~(v3) keeps the search algorithm of RIsearch2 unchanged: targets are concatenated and indexed in a suffix array, then each query is scanned against the index for near-complementary seeds
The seeds are then extended in both directions with Gotoh's three-state affine-gap recurrence~\citep{gotoh1982improved}, scoring stacked pairs from a nearest-neighbour dinucleotide model that penalising bulges and internal loops.
The implementation, by contrast, was rewritten from scratch in Rust.
Most of the code is safe Rust, which guarantees memory safety; unchecked memory access is confined to the performance-critical inner loops of seed search and extension and to the zero-copy loading of the on-disk index.
The suffix array is now built by libsais~\citep{libsais}, a linear-time construction that can run in parallel across cores, in place of the single-threaded libdivsufsort~\citep{libdivsufsort} used by RIsearch2.
The energy parameter tables were regenerated from ViennaRNA duplex energies, so calculated hybridisation energies differ from those of RIsearch2.

\subsection*{Seamless installation}

RIsearch2 is distributed as C source code that each user must compile by hand, which requires a C toolchain and leaves build settings and dependency versions to each machine.
The new version needs no manual build: it installs with a single command from the Python and Rust package registries, \href{https://pypi.org/project/risearch/}{PyPI} and \href{https://crates.io/crates/risearch}{Crates.io}, or as a prebuilt binary from each \href{https://github.com/saiden89/risearch/releases}{GitHub release}.

\subsection*{Improved usability}

While RIsearch2 can return many overlapping alignments for one duplex, RIsearch~(v3) by default reports one alignment per duplex: seed extensions that end on exactly the same query and target coordinates are collapsed to the one with the lowest (most favorable) energy.
For example, two seeds at neighbouring positions of an siRNA that extend to the same duplex yield one row instead of two near-identical ones, unless collapsing is disabled.
RIsearch~(v3) also validates its inputs more strictly than RIsearch2, rejecting malformed input with an error that names the offending record or line.
The command line is rebuilt around task-oriented subcommands, while legacy RIsearch2 flags are still accepted to ease migration of existing scripts.

\subsection*{Enhanced interoperability}

RIsearch2 was designed as a standalone command-line tool, so embedding it in a Python analysis meant spawning a subprocess and parsing its text output.
The new version of RIsearch instead ships PyO3 bindings in the same wheel: \texttt{risearch.index()} builds an on-disk index, \texttt{risearch.TargetRegistry.open()} loads it as a reusable memory-mapped handle, and \texttt{risearch.search()} returns hits as a Polars DataFrame through the Arrow C stream interface, with no intermediate files and no text parsing:

\begin{verbatim}
import risearch
risearch.index("targets.fa", "index/")
store = risearch.TargetRegistry.open("index/")
df = risearch.search(queries, store, matrix="t04")
\end{verbatim}

Searches run without blocking other Python threads in the same process, and the energy parameter set can be chosen per call.

\subsection*{Extensible energy parameters and engine}

The energy parameters are stored as tables separate from the search code: the bundled sets are generated from ViennaRNA by a script in the RIsearch repository, and users can load their own table at run time, so new thermodynamic parameters require no change to the search engine and no knowledge of Rust internals.
The engine itself is modular, and its behaviour is verified against independent reference implementations through property-based and metamorphic tests, complemented by mutation testing and bounded model-checking proofs (Kani), so contributors can extend itwith confidence.

\subsection*{Faster execution time}

On an AMD EPYC~9554 server, with both tools searching the same 63 siRNAs against the same chromosome-X transcript set, RIsearch~(v3) is $2$--$3\times$ faster than RIsearch2 (Supplementary Figures~\ref*{supp-fig:risearch-comparison}--\ref*{supp-fig:risearch-compression-scaling}; Supplementary Methods, ``RIsearch core search benchmark'').
The suffix array creation rewrite also yielded about $2\times$ faster execution times (Supplementary Figure~\ref*{supp-fig:indexing}; Supplementary Methods, ``Index construction'').
Because on-disk index is now memory-mapped and read without deserialisation, multi-gigabyte genome indexes are  available immediately and shared across processes by the operating system, rather than read into memory on every run.

% JG: "Our siOFF (siRNA off-target assessment pipeline) is a complete reimplementation of the siRNA off-target pipeline introduced by \citet{alkan2017risearch2}"; "porting" -> "adapting".
% JG: the "three practical barriers" sentence repeats the introduction; here the text must be more elaborate than a repeat.
% JG: point to precisely where in the Supplementary Methods.
\section{siOFF: A reimplemented siRNA Off-Target assessment pipeline}\label{sec:pipeline}
Alongside the new version of RIsearch, we distribute \textbf{siOFF}, a complete reimplementation of the siRNA off-target pipeline of \citet{alkan2017risearch2} (\cref{fig:summary}).
Beyond adapting the workflow from Python~2.7 to Python~3, the rewrite targets three practical barriers that limited adoption of the original tool: installation, ease of use, and integration into existing workflows.
The underlying four-stage thermodynamic model---search, annotation intersection, accessibility, and a per-siRNA Boltzmann partition function---is preserved unchanged and described in full in \ref*{supp-sec:pipeline-stages}.


\subsection*{One-command installation}
The original pipeline requires a Python~2.7 interpreter, a separately compiled RIsearch2 C binary, and installation of external software \texttt{RNAplfold} (ViennaRNA) and \texttt{bedtools}~\citep{quinlan2010bedtools}.
siOFF removes the need for external dependency installation: a single \texttt{pip install sioff-rna} brings in the RIsearch software as a pre-compiled wheel, the ViennaRNA~2.7.2 Python API~\citep{lorenz2011vienna} in place of the \texttt{RNAplfold} executable, and an in-memory NCLS interval index~\citep{ncls} in place of \texttt{bedtools}.
All transitive dependencies are resolved and pinned through a universal \texttt{uv} lockfile, so installations are reproducible and free of version conflicts across machines.

\subsection*{Easier to use}
The legacy pipeline was a single \texttt{argparse} script driven by a fixed set of command-line flags.
siOFF retains backwards-compatible flags but adds a YAML configuration mode: every option---inputs, the RIsearch energy model, seed geometry, and energy threshold, the saturation parameters $\alpha$, $\gamma$, and $\theta$, and the output format---can be declared in one file, validated at startup against a typed OmegaConf schema.
A single configuration file captures an entire run, making analyses straightforward to reproduce and to share between collaborators.
For large screens, siOFF integrates natively with Slurm: one orchestrator submits the indexing, accessibility, and off-target stages as dependent batch jobs with per-stage resource requests, and a single configuration can fan one screen out across multiple transcriptomes as independent job groups, replacing hand-written submission scripts.

A single temperature option (default $37\,^\circ$C) has now been implemented. The temperature propagates to both the ViennaRNA accessibility calculation and the Boltzmann partition function, so that hybridisation and opening energies are always evaluated on a consistent thermodynamic scale---directly replicating the intent of the original RIsearch model~\citep{wenzel2012risearch} while exposing the folding condition as a first-class parameter (Supplementary Methods \ref*{supp-sec:energy-params}).

\subsection*{Easy integration}
Because the legacy pipeline communicated through intermediate files and subprocesses, embedding it in a larger analysis meant shelling out to a script and parsing its text output.
siOFF instead exposes its stages as an importable Python API: the search, intersection, accessibility, and probability services can each be called directly and return native Polars DataFrames, so a user can drop any stage into an existing Python workflow and pass results to other Python-native tools in a single process, with no intermediate files and no text parsing.
Based on RIsearch software bindings (\cref{subsec:pybindings}), siOFF thus serves both as a turnkey command-line tool and as a library component of off-target ad hoc analyses.

% JG: point to precisely where in the Supplementary Methods (both references).
\subsection*{Enhanced execution time}
The pipeline was benchmarked on an AMD EPYC 9554 (64 cores, 32 workers) using subsets of the genome-wide siRNA efficacy library of \citet{huesken2005design} searched against the GRCh38/hg38 reference genome~\citep{schneider2017grch38} and intersected against an H1299 RNA-seq expression profile (dataset, genome indexing, expression annotation, and benchmark protocol detailed in \ref*{supp-sec:benchmark-protocol}), siOFF completes the intersection, accessibility annotation, probability computation, and output stages $7.7{-}10.7\times$ faster than the legacy Python~2.7 implementation across batch sizes of 500--2{,}000 siRNAs (five randomised trials per batch size; Supplementary Figure~\ref*{supp-fig:pipeline}).
Peak memory plateaus at approximately 7\,GB regardless of batch size, whereas the legacy pipeline grows from $\sim$8.5\,GB to $\sim$10\,GB as batch size increases (Supplementary Figure  \ref*{supp-fig:pipeline}).

Separately, the one-time genome-wide accessibility pre-computation---folding the entire hg38 reference at $W{=}80$, $L{=}40$, $u{=}30$ with 64 workers---runs $\sim$34\% faster than the legacy \texttt{RNAplfold} script (8\,h\,46\,m vs.\ 13\,h\,12\,m) and writes profiles $\sim$42\% smaller on disk (111.8\,GB vs.\ 192.6\,GB), at the cost of higher but still modest peak memory ($\sim$3.6\,GB vs.\ $\sim$309\,MB; Supplementary Figure~\ref*{supp-fig:accessibility} and Supplementary Table~\ref*{supp-tab:acc-bench}).

\section{Final Remarks}

Here, we have exploited the emerging possibilities of new software ecosystems to reimplement RIsearch and the siOFF pipeline, with enhancements in installation, maintenance, integration with other tools and speed.
Both tools can be used from the command line and within Python workflows by anyone screening RNA interactions or siRNA off-targets at transcriptome scale.
Future work will extend siOFF to the combined off-targeting potential of pooled siRNAs and integrate RIsearch with further analysis frameworks.


\section{Author contributions statement}
S.R. reimplemented RIsearch2 in Rust and provided the framework for the siOFF pipeline.
L.F. designed and implemented the siOFF pipeline, ran all benchmarks, and drafted the manuscript.
C.A. designed the temperature option and tested the software.
J.G. supervised the project.
All authors commented on, contributed to, and approved the final manuscript.

\section{Funding}
This work is supported by the Novo Nordisk Foundation (NNF24OC0087112).

\section{Conflict of interest}
None declared.

% JG: is the benchmark data the same as in RIsearch2? If so we should not redistribute it.
\section{Code and data availability}
%TODO: update when transfer ownership happens.
The source code for RIsearch and siOFF is available on GitHub (\faGithub~\href{https://github.com/saiden89/risearch}{saiden89/risearch}, \faGithub~\href{https://github.com/lorenzo-22/siOFF}{lorenzo-22/siOFF}), and the packages on PyPI (\faPython~\href{https://pypi.org/project/risearch/}{risearch}, \faPython~\href{https://pypi.org/project/sioff-rna/}{sioff-rna}) and Crates.io (\faRust~\href{https://crates.io/crates/risearch}{risearch}).
All distribution channels are summarised at \faGlobe~\href{http://rth.dk/resources/risearch/}{rth.dk/resources/risearch}.
% RIsearch and siOFF source code is available at [GitHub URL]. Benchmark datasets and scripts
% used to reproduce the results in this paper are deposited at [Zenodo DOI].

\bibliographystyle{oup-abbrvnat}
\bibliography{reference}

% Supplementary material is a separate document: see supplementary.tex.

\end{document}
